\section{Challenges, Consultation with TA, and Week 04 Roadmap}
\label{section:roadmap}

\subsection{Technical Challenges Encountered in Week 03}
During our Week 03 exploration, our team identified several technical hurdles:
\begin{itemize}
    \item \textbf{a. Channel capacity inflation in single-stage models:} While Shift-ResMLP successfully achieved 72.92\% accuracy with parameter-free spatial mixing, its isotropic single-scale nature required a large channel width ($D = 384$) and 12 blocks, accumulating 14.26 million parameters. In contrast, multi-stage pyramids achieve comparable performance with 1.79 million parameters, confirming that isotropic backbones suffer from intrinsic capacity inefficiency.
    \item \textbf{b. Regularization requirements at higher resolutions:} On Tiny-ImageNet ($64*64$), our base pyramid model exhibited overfitting due to increased spatial token counts. Introducing stochastic depth (DropPath = 0.05) and heavier weight decay was essential to unlock strong generalization (+5.08\% improvement to 62.93\%), indicating that regularization becomes increasingly critical as resolution scales up.
    \item \textbf{c. Disparate peer training pipelines:} Because team members worked on diverse compute environments (local DDP clusters, single CUDA GPUs, and Kaggle T4 accelerators), differing batch sizes, learning rates, and augmentation suites made cross-model comparison nuanced. A primary goal for Week 04 is unifying all top-performing candidates under an identical benchmarking harness.
\end{itemize}

\subsection{Questions for TA Consultation}
We have prepared two technical research questions for our TA:

\begin{quote}
\textbf{Q1: Hardware Inference Profiling and Re-parameterization Verification} \\
Pyramid-ResMLP demonstrates high parameter efficiency (1.94M training parameters fusing into 1.79M deployment parameters via RepTokenMix) and competitive accuracy (72.18\% on CIFAR-100). To conclude our research in Week 04, we intend to empirically validate that structural re-parameterization delivers genuine zero-latency overhead in real deployment. What benchmarking protocol does the TA recommend for profiling inference latency and throughput (e.g., measuring TensorRT/CUDA latencies across batch sizes $b \in \{1, 32, 128\}$, GPU memory, and kernel launch overhead) to conclusively verify the deployment advantage of the fused linear operator over its unfused multi-branch counterpart?
\end{quote}

\begin{quote}
\textbf{Q2: Stage Depth Allocation and Regularization for High-Resolution Scaling} \\
On Tiny-ImageNet ($64 \times 64$), introducing stochastic depth ($\text{DropPath} = 0.05$) and expanding channel widths to $[112, 224, 448]$ boosted Pyramid-ResMLP from 57.85\% to 62.93\% (+5.08\%). To extend Pyramid-ResMLP toward higher-resolution vision benchmarks in Week 04, should our team prioritize asymmetric stage depth scaling (such as concentrating depth in Stage 3 with a $[2, 2, 6, 2]$ ratio following ConvNeXt and PoolFormer) or balanced channel expansion? Furthermore, to maximize generalization without large-scale pretraining, would the TA advise exploring Teacher-Student knowledge distillation (e.g., from a compact pretrained CNN teacher) or continuing to refine data augmentations and stochastic depth?
\end{quote}

\subsection{Research Roadmap for Week 04}
To conclude our project milestone, our team plans the following research activities:
\begin{itemize}
    \item \textbf{a. Comprehensive hardware inference profiling:} Implement a rigorous benchmarking pipeline measuring inference latency (ms/batch), throughput (images/sec), and GPU memory footprint using CUDA events with warmup and TensorRT/ONNX Runtime to empirically validate the zero-overhead property of fused RepTokenMix.
    \item \textbf{b. Macro-architecture scaling on Tiny-ImageNet:} Systematically evaluate asymmetric stage depth allocations and capacity configurations for Pyramid-ResMLP at $64 \times 64$ resolution.
    \item \textbf{c. Final ablation suite and unified benchmarking:} Conduct complete ablation studies on downsampling stem kernels, RepTokenMix branch contributions, and regularization schedules under unified seeds to conclude the final technical report.
\end{itemize}
